(1) 设 \(\Phi_0=h/(2e)=2\pi\hbar/(2e)\)，总并联超电流为
\[
I=I_1\sin(\delta+\phi/2)+I_2\sin(\delta-\phi/2)
=(I_1+I_2)\cos(\phi/2)\sin\delta
+(I_1-I_2)\sin(\phi/2)\cos\delta.
\]
记复振幅
\(Z(\phi)=I_1e^{i\phi/2}+I_2e^{-i\phi/2}=I_c(\phi)e^{i\alpha(\phi)}\)，则
\(I=I_c(\phi)\sin[\delta+\alpha(\phi)]\)。
Josephson势能为
\[
U_J=-E_1\cos(\delta+\phi/2)-E_2\cos(\delta-\phi/2)
=-E_{\rm eff}(\phi)\cos[\delta+\alpha(\phi)],
\quad E_{\rm eff}=\frac{\hbar I_c}{2e}.
\]
此处的有效单结是两弱结在共同相位 \(\delta\) 上干涉后的等效势，不是单弱结电流的代数求和。

(2) 对共同相位取绝对最大值得
\[
I_c(\phi)=|Z|=\sqrt{I_1^2+I_2^2+2I_1I_2\cos\phi}
=\sqrt{(I_1+I_2)^2\cos^2\frac\phi2+
(I_1-I_2)^2\sin^2\frac\phi2}.
\]
在整数磁通处 \(I_c=I_1+I_2\)，在半整数磁通处
\(I_c=|I_1-I_2|\)；等结 \(I_1=I_2=I_0\) 时
\(I_c=2I_0|\cos(\pi\Phi_{\rm ext}/\Phi_0)|\) 有精确节点，而不等结保留非零极小值。平方根内只含 \(\cos\phi\)，故所有可观测临界电流以 \(\Phi_0\) 为周期；\(\alpha\) 变化时可同时平移共同相位，避免误判 \(2\Phi_0\) 周期。

(3) 偏置源作功的有效势为
\[
U(\delta)=-E_{\rm eff}\cos(\delta+\alpha)
-\frac{\hbar I_b}{2e}\,\delta.
\]
设 \(i=|I_b|/I_c\in[0,1)\)，作平移 \(\varphi=\delta+\alpha\)（若负偏置相应反向选鞍点），极值满足
\(\sin\varphi=I_b/I_c\)。
对于正偏置，一极小值在 \(\varphi_m=\arcsin i+2\pi n\)，沿逃逸方向的相邻鞍点在 \(\varphi_s=\pi-\arcsin i+2\pi n\)，故
\[
\Delta U=2E_{\rm eff}\left[\sqrt{1-i^2}-i\arccos i\right].
\]
若 \(i\to1^-\)，令 \(\epsilon=1-i\)，利用
\(\sqrt{1-i^2}-i\arccos i=(2\sqrt2/3)\epsilon^{3/2}+O(\epsilon^{5/2})\)，从而
\[
\Delta U=\frac{4\sqrt2}{3}E_{\rm eff}(1-i)^{3/2}
+O((1-i)^{5/2}).
\]
若 \(I_c=0\)（等结恰半磁通），无稳态束缚势和上述除以 \(I_c\) 的公式；若忽略电容，只是省略惯性和逃逸速率的动力学影响，静态势垒形状仍成立。

(4) 零偏置下对 \(\delta\) 极小化有 \(F(\Phi_{\rm ext})=-E_{\rm eff}(\Phi_{\rm ext})\)，故
\[
I_{\rm circ}=-\frac{\partial F}{\partial\Phi_{\rm ext}}
=\frac{\partial E_{\rm eff}}{\partial\Phi_{\rm ext}}
=-\frac{I_1I_2\sin\phi}{I_c(\phi)}
\quad (I_c>0).
\]
在等结 \(I_1=I_2=I_0\) 且未跨节点的区域，公式等于
\(-I_0\sin(\phi/2)\operatorname{sgn}[\cos(\phi/2)]\)；半磁通的基态出现简并尖点，导数从两侧不同，不能说该点有唯一平衡环流。若 \(I_2\to0\)，\(I_c\to I_1\) 而环流趋零，符合一个弱结构不成双路径干涉。有限自感 \(L\) 时穿环总磁通为 \(\Phi_{\rm ext}+LI_{\rm circ}\)（符号依环流方向约定），需自洽求解相位，以上简单 \(\Phi_0\) 周期临界曲线和尖点会受屏蔽修正。
